\section{Manual (Plataforma web) - Apartado de gestión de reservas e instalaciones}

Al ingresar a la página de reservas e instalaciones, se visualiza un listado con todas las instalaciones disponibles en el club[cite: 16]. De cada una de ellas se puede observar su nombre, tipo, capacidad máxima y valor asignado[cite: 16].

\subsection{Detalle de instalación y gestión de reservas}
Al entrar a una instalación en particular, la plataforma muestra una descripción más extensa del espacio[cite: 16]. Desde esta sección es posible visualizar tanto las reservas activas (por ejemplo, del socio 1000) como las reservas históricas[cite: 16].

Además, el sistema permite crear de manera manual una reserva para un socio específico[cite: 16]. 

Pasos para registrar una reserva[cite: 16]:
\begin{itemize}
    \item Ingresar el número de socio (por ejemplo, socio con número 1001)[cite: 16].
    \item Seleccionar la fecha y el turno deseado (por ejemplo, 19 de agosto de 2026 a las 11 de la mañana)[cite: 16].
    \item Al registrarla, la reserva quedará en estado pendiente de pago[cite: 16].
    \item Esta acción se reflejará automáticamente en el perfil del socio titular, donde figurará que posee una reserva con pago pendiente[cite: 16].
\end{itemize}

\subsection{Creación de nuevas instalaciones}
Desde este apartado también se pueden dar de alta nuevos espacios físicos para el club[cite: 16]. Para la creación de una instalación, se deben configurar los siguientes parámetros[cite: 16]:

\begin{itemize}
    \item \textbf{Nombre y tipo:} Denominación y categoría del espacio (por ejemplo, Parrilla 4 de tipo social)[cite: 16].
    \item \textbf{Capacidad máxima:} Cantidad límite de personas permitidas (por ejemplo, 100 personas)[cite: 16].
    \item \textbf{Valor por turno:} Costo de la reserva (por ejemplo, 10 mil pesos)[cite: 16].
    \item \textbf{Duración del turno:} Tiempo de uso asignado (por ejemplo, dos horas)[cite: 16].
    \item \textbf{Tiempo mínimo de cancelación:} Plazo límite establecido para cancelar la reserva[cite: 16].
\end{itemize}

Una vez completados todos los datos, la nueva instalación quedará creada exitosamente en el sistema[cite: 16].
